%===============================================================================
% ファイル名: master_ja_sample.tex
% 作成者：サイトウ リョウタ（2024年度山上研究室 修士修了）
% 作成日：2025年2月19日
% 
% これは京大土木系の修士論文作成のためのサンプルLaTeXファイルです．
% 「コンパイル環境をいじらないで使える！」を目指してま～す☆
% LuaLaTeXでコンパイルするとバッチリだよ！
%
% ちなみに，すいちゃんは～？今日もかわいい～！
%===============================================================================
\documentclass[a4paper, 11pt]{report}
\usepackage{luatexja}

% Required packages
\usepackage{amsmath, amsthm, amssymb, amsfonts}
\usepackage{graphicx}
\usepackage{physics}
\usepackage{subcaption}
\usepackage{geometry}
\usepackage{setspace}
\usepackage{hyperref}
\usepackage{titlesec}
\usepackage{fancyhdr}
\usepackage{bookmark}
\usepackage{caption}
\usepackage{enumitem}
\usepackage{siunitx}
\usepackage{luatexja-fontspec}
\usepackage{tocloft}
\usepackage{indentfirst} % Enable indentation for the first paragraph after a chapter or section
\usepackage{cite} % ref
\usepackage{float} % for \begin{figure}[H] option in figures, fix in the position
\usepackage{xcolor}
\usepackage{ifthen}
\usepackage{bm}
\usepackage{booktabs}
\usepackage[ppl]{mathcomp}
\usepackage{multirow}

\makeatletter
\renewcommand{\@cite}[2]{\textsuperscript{#1)}} % "1)"

\newcommand{\supercite}[2]{\cite{#1}\textsuperscript{,}\cite{#2}}

% Reference style
\renewenvironment{thebibliography}[1]
{\section*{\fontsize{18pt}{22pt}\selectfont 参考文献}
    \list{\arabic{enumi})}
        {\settowidth\labelwidth{#1)}%
        \leftmargin\labelwidth
        \advance\leftmargin\labelsep
        \usecounter{enumi}%
        \let\p@enumi\@empty
        \renewcommand\theenumi{\arabic{enumi})}}%
    \sloppy\clubpenalty4000\widowpenalty4000\sfcode`\.\@m}
{\endlist}
\urlstyle{same}   

% Set font to MS Mincho
\setmainjfont{MS Mincho}[
    RawFeature={+palt, +kern},
    LetterSpace=12.5 % Increase spacing between characters
]
\setsansjfont{MS Gothic}[
    RawFeature={+palt, +kern},
    LetterSpace=10.0
]
\renewcommand{\bfseries}{\gtfamily}

% Set font to Times New Roman for English
\setmainfont{Times New Roman}[
    RawFeature={+palt, +kern},
    LetterSpace=2.0
]
\setsansfont{Arial}[
    RawFeature={+palt, +kern},
    LetterSpace=1.5
]
\setmonofont{Courier New}[
    RawFeature={+palt, +kern},
    LetterSpace=1.0
]

% Page margins settings
\geometry{
    left=20mm,      % Left margin
    right=20mm,     % Right margin
    top=25mm,       % Top margin
    bottom=25mm     % Bottom margin
}

% Line spacing
\setstretch{1.2}
\renewcommand{\baselinestretch}{1.2}
\linespread{2.5}
\newcommand{\HRule}[1]{\rule{1\linewidth}{#1}}

% Rename figure and table captions to Japanese
\renewcommand{\figurename}{\normalfont\gtfamily 図} % Figure label to Japanese
\renewcommand{\tablename}{\normalfont\gtfamily 表} % Table label to Japanese
\renewcommand{\thefigure}{\thechapter-\arabic{figure}}
\renewcommand{\thetable}{\thechapter-\arabic{table}}
\newcommand{\figureref}[1]{図\ref{#1}}
\newcommand{\tableref}[1]{表\ref{#1}}

% Page numbering
\pagestyle{fancy}
\fancyhf{}
\fancyfoot[C]{\thepage} % Page number centered in the footer
\renewcommand{\headrulewidth}{0pt}

% =====================================================
% REMOVE FOR SUBMISSION!!!
% Set color for my eyes
% Disable or delete the following settings to revert 
% the background and text colors to normal settings!!!
% 
% Conditional flag: Submission mode (1: Submission, 0: Debug mode)
\newcommand{\submissionmode}{1} % Set to 1 for submission
% 
% =====================================================

% NOTE: REMOVE FOR SUBMISSION!
\newcommand{\debugnote}{
    \ifthenelse{\submissionmode=0}{
        \textbf{\textcolor{red}{\LARGE{REMOVE THIS COLOR SETTING FOR SUBMISSION!}}} \\
    }{}
}

\ifthenelse{\submissionmode=0}{
    \pagecolor{black}
    \color{white}
}{}

% =====================================================
% =====================================================
\titleformat{\chapter}
    {\normalfont\fontsize{20pt}{24pt}\gtfamily} % Style and font size
    {第\hspace{2pt}\thechapter 章} % Chapter format: "\u7b2c" ("第") and "\u7ae0" ("章") for "Chapter"
    {20pt} % Space between chapter number and title
    {\fontsize{20pt}{10pt}\selectfont} % Font size for title text

% Section and subsection title formatting
\titleformat{\section}
    {\normalfont\fontsize{12pt}{12pt}\gtfamily} % Style and font size
    {\thesection}{1em}{}

\titleformat{\subsection}
    {\normalfont\fontsize{12pt}{12pt}\gtfamily} % Style and font size
    {\thesubsection}{1em}{}

% Contents title to Japanese
\renewcommand{\contentsname}{\normalfont\gtfamily 目次} % Set "\u76ee\u6b21" ("目\u6b21") for "Contents"

% number of chaper, section, subsection
\renewcommand{\thechapter}{\arabic{chapter}}
\renewcommand{\thesection}{\arabic{chapter}.\arabic{section}}
\renewcommand{\thesubsection}{\thesection.\arabic{subsection}}

% tocloft
\renewcommand{\cftchappresnum}{第}
\renewcommand{\cftchapaftersnum}{章\hspace{0.5em}}

% for section and subsection
\renewcommand{\cftsecpresnum}{} 
\renewcommand{\cftsecaftersnum}{\hspace{0.5em}}
\renewcommand{\cftsubsecpresnum}{}
\renewcommand{\cftsubsecaftersnum}{\hspace{0.5em}}

% to output % (percent)
\newcommand{\percent}{\%}

% Adjust spacing in the table of contents
\setlength{\cftchapnumwidth}{4.0em} % Width for "第1章"

\begin{document}
\clearpage
\pagenumbering{roman} % i, ii, iii
\setlength{\baselineskip}{19.5pt}

% Title page
\thispagestyle{empty}

% Logo of Kyoto University
\noindent
\begin{minipage}[t]{\linewidth}
    \vspace{3.5cm}
    \raggedleft % Align to the right
    \raisebox{-1cm}[0pt][0pt]{
    \includegraphics[width=0.18\linewidth]{./img/title/kyodai.png}
    }
\end{minipage}

\vspace{-1.5cm}
\noindent
{\fontsize{12pt}{12pt}\gtfamily
\hspace{-0.3em}
\setlength{\baselineskip}{12pt}
\begin{minipage}[t]{\linewidth}
    京都大学大学院工学研究科 \\[-1.3em]
    社会基盤工学専攻修士論文 \\[-1.3em]
    令和７年２月 \\[-1.3em]
\end{minipage}
}

\noindent
{\fontsize{12pt}{12pt}\gtfamily
\hspace{-0.3em}
\setlength{\baselineskip}{12pt}
\begin{minipage}[t]{\linewidth}
    Master's Thesis \\[-1.3em]
    Department of Civil and Earth Resources Engineering \\[-1.3em]
    Graduate School of Engineering \\[-1.3em]
    Kyoto University \\[-1.3em]
    February 2025 \\[-1.7em]
\end{minipage}
}
\HRule{1.0pt}
\begin{center}
    \vspace{2.5cm}
    \debugnote
    {\fontsize{18pt}{36pt}\selectfont 
    \LaTeX による修士論文作成フォーマット
    } \\
    \vspace{6.5cm}
    {\fontsize{16pt}{18pt}\selectfont 
    京都大学大学院 \hspace{1em}工学研究科 \hspace{1em}社会基盤工学専攻 \\ 
    アイドル講座\hspace{1em}バーチャル音響学分野 \\ 
    星街\hspace{1em}すいせい \\
    }
\end{center}
\newpage

% Abstract
\begin{center}
    {\fontsize{12pt}{26pt}\selectfont \gtfamily 論文要旨} % 20pt font size, 24pt line spacing, gothic
\end{center}
\par \indent
ここから要旨を書き始めます．
「論文要旨」はセンタリングされます．

改行する際にはこのように空行を一行入れます．
「彗星のごとく現れたスターの原石！バーチャルアイドルの星街すいせいでーす！」
「彗星のごとく現れたスターの原石！バーチャルアイドルの星街すいせいでーす！」
「彗星のごとく現れたスターの原石！バーチャルアイドルの星街すいせいでーす！」
「彗星のごとく現れたスターの原石！バーチャルアイドルの星街すいせいでーす！」
「彗星のごとく現れたスターの原石！バーチャルアイドルの星街すいせいでーす！」
このようにすれば改行されません．
\newpage

% Table of contents
\tableofcontents
\newpage

\clearpage
\pagenumbering{arabic}

% Main content starts here
% ------------------------------------------------------------------------------
% ---chapter 1.-----------------------------------------------------------------
% ------------------------------------------------------------------------------
\chapter{はじめに}
この\LaTeX フォーマットは，京都大学大学院工学研究科社会基盤工学専攻・都市社会工学専攻の修士論文を執筆するための基本的な
機能付きのテキストファイルです．
とりあえず，2024年度の修士論文フォーマットには対応しています．

使用する\LaTeX コンパイラはLuaLaTeX環境を想定しています．
pLaTeX環境等をお使いの方は，適宜プリアンプルの修正を行ってください．

また，本フォーマットにはデバッグモードを実装しています．
出力されるpdfがブラック背景に白文字になるので，目の疲れを低減することができます．
submissionmodeの1を0に変更すればデバッグモードになります．

\section{新しい節}
節はsectionコマンドで入れます．

\subsection{新しい項}
項はsubsectionコマンドを用います．

% ------------------------------------------------------------------------------
% ---chapter 2.-----------------------------------------------------------------
% ------------------------------------------------------------------------------

\chapter{論文の執筆}
\section{図の挿入}
\begin{figure}[H]
    \centering
    \includegraphics[width=0.8\linewidth]{./img/chap2/computer_tokui_boy.png}
    \caption{コンピューターの得意な男の子}
    \label{fig_pcboy}
\end{figure}
[H]を使うと.texの入力通りの場所にfloatを配置できますが，文書が崩れる可能性があります．
基本的には[p]または[t]を用いるのがおすすめです．
このサンプルでは文章量の都合上[H]指定しています．
そのほかの設定やsubfigureの環境については調べてみてください．

\section{表の挿入}
\begin{table}[H]
    \centering
    \renewcommand{\arraystretch}{.8}
    \setlength{\tabcolsep}{12pt}
    \caption{星街すいせいのプロフィール\cite{11}}
    \begin{tabular}{c|c}
        \toprule
        \textbf{項目} & \textbf{内容} \\ 
        \midrule
        誕生日 & 3月22日 \\ 
        デビュー日 & 2018年3月22日 \\ 
        身長 & \SI{160}{cm} \\ 
        ユニット & ホロライブ0期生 \\ 
        メイク担当 & 手島nari \\ 
        ファンネーム & 星詠み \\ 
        配信タグ & \#ほしまちすたじお \\ 
        ファンアートタグ & \#ほしまちぎゃらりー \\ 
        \bottomrule
    \end{tabular}
    \label{tab_suisei_profile}
\end{table}
このように表を挿入することができます．
なぜか[H]を指定すると行間が大きくなってしまうので0.8を指定しましたが，[t]なら1.8程度でいいと思います．

\section{数式の挿入}
数式は文中への挿入またはequation環境での挿入の二種類に分かれます．
文中に挿入する際は，\(L = 4\pi R^2 \sigma T^4\)このように追加します．

独立表示させるには，
\begin{equation}
    r = \frac{a(1 - e^2)}{1 + e \cos\theta}
\end{equation}
このように書きます．

\section{引用・参照}
\tableref{tab_suisei_profile}において引用を行いましたが，引用はciteコマンドを用いています．
また，表や図，数式の参照を行う際はref系のコマンドで行います．

引用については，thebibliography環境を用いています．
自分が執筆した際は手動入力を行いました．
bibtex環境を用いてもいいと思いますが，コンパイルの手順や必要packageについてはご確認ください．

\section{その他}
SI単位系にはSIコマンドを用いています．
\SI{10}{m/s}等入力して，SI単位系付きの数値を適切に入力できます．

% ------------------------------------------------------------------------------
% ---chapter 3.-----------------------------------------------------------------
% ------------------------------------------------------------------------------
% chapter 3以降も同様に追加していきます．
% \chapter{以降の章タイトル}

% ------------------------------------------------------------------------------
% Acknowledgement
% ------------------------------------------------------------------------------

\newpage
\addtocontents{toc}{\vspace{3em}} % vspace in TOC
\section*{\fontsize{18pt}{22pt}\selectfont 謝辞} % Acknowledgement in TOC
\addcontentsline{toc}{section}{謝辞}
ここから謝辞が入ります．（主査の教授\rightarrow 副査の教授 \rightarrow 副査の准教授の先生の順，その後に
お世話になった先生や技術職員の方，他の学生等について言及する）

% ------------------------------------------------------------------------------
% Reference and Cited Works
% ------------------------------------------------------------------------------

\newpage
\addtocontents{toc}{\vspace{2em}} % vspace in TOC
\addcontentsline{toc}{section}{参考文献} % Reference in TOC
\begin{thebibliography}{99}
    % chapter 1
    \bibitem{11} 星街すいせい，ホロライブ公式サイト，\url{https://hololive.hololivepro.com/talents/hoshimachi-suisei/}，（参照，2025-2-19）.
    % chapter 2
    % chapterごとに分けておくと後で楽，pbibtexでのコンパイルを組み合わせてもいいと思います．
\end{thebibliography}

% ------------------------------------------------------------------------------

\end{document}

% ------------------------------------------------------------------------------
% ---end document---------------------------------------------------------------
% ------------------------------------------------------------------------------




